\subsection{Expectation--maximization (EM)}
\label{subsec:expectation-maximization}

Many statistical models are easiest to describe by introducing a latent or missing variable.
This is the probabilistic version of augmentation: the auxiliary variable $z$ makes the complete-data model simple even though the observed-data likelihood is harder.
Suppose the observed data are $x$, the latent data are $z$, and the model is specified by a joint density or mass function $p_\theta(x,z)$.
If $z$ were observed, then we could often optimize the complete-data model directly.
The difficulty is that maximum likelihood depends on the marginalized quantity
\[
  p_\theta(x) = \int p_\theta(x,z)\,dz,
  \qquad
  \ell(\theta;x) \coloneqq \log p_\theta(x),
\]
with the integral replaced by a sum when $z$ is discrete.
The awkward step is that the logarithm sits outside the integral, so even a simple complete-data model can lead to a difficult observed-data optimization problem.
EM replaces this direct maximization by alternating two simpler steps.
With the current iterate \(\theta^{(t)}\), the E-step computes the current conditional law
\[
  q_t(z)\coloneqq p_{\theta^{(t)}}(z\mid x),
\]
and the M-step maximizes the averaged complete-data log-likelihood
\[
  Q(\theta\mid \theta^{(t)})
  \coloneqq
  \E_{Z\sim q_t}\!\left[\log p_\theta(x,Z)\right].
\]
Later in the section we will rewrite this same alternation as coordinate ascent on a free-energy functional, which is the bridge to variational inference.
Classical examples include Gaussian mixture models, hidden Markov models, and Gaussian models with missing data.

\subsubsection{Free energy, the \texorpdfstring{$Q$}{Q}-function, and coordinate ascent}
\label{subsubsec:em-q-function}

Let $\theta^{(t)}$ be the current iterate and define the current posterior law of the latent variable by
\[
  q_t(z) \coloneqq p_{\theta^{(t)}}(z\mid x).
\]
The associated EM surrogate is
\begin{equation}
  Q(\theta\mid \theta^{(t)})
  \coloneqq
  \E\bigl[\log p_\theta(x,Z)\mid X=x,\theta^{(t)}\bigr]
  =
  \int \log p_\theta(x,z)\,q_t(z)\,dz.
  \label{eq:em-q-function}
\end{equation}
The E-step computes the conditional law $q_t$, and the M-step maximizes the surrogate \eqref{eq:em-q-function} with respect to $\theta$.
$Q(\theta\mid \theta^{(t)})$ averages the complete-data log-likelihood under the current belief about the missing data, so the M-step often reduces to a familiar weighted maximum-likelihood problem.
The E-step updates the auxiliary latent law, and the M-step optimizes with that law held fixed.

\begin{algorithm}
  \caption{Expectation--maximization (EM)}
  \label{alg:em}
  \begin{algorithmic}[1]
    \Require Observed data $x$, joint model $p_\theta(x,z)$, initial guess $\theta^{(0)}$, stopping tolerance $\varepsilon$
    \Ensure Approximate maximizer of the observed-data log-likelihood
    \For{$t=0,1,2,\dots$}
      \State \textbf{E-step:} compute the conditional law $q_t(z)=p_{\theta^{(t)}}(z\mid x)$
      \State Form $Q(\theta\mid \theta^{(t)}) = \E_{Z\sim q_t}\!\left[\log p_\theta(x,Z)\right]$
      \State \textbf{M-step:} choose $\theta^{(t+1)} \in \argmax_{\theta} Q(\theta\mid \theta^{(t)})$
      \If{$\norm{\theta^{(t+1)}-\theta^{(t)}}_2 \le \varepsilon$}
        \State \Return $\theta^{(t+1)}$
      \EndIf
    \EndFor
  \end{algorithmic}
\end{algorithm}

The M-step does not have to be solved exactly to preserve the usual ascent property.
Any update that increases $Q(\theta\mid \theta^{(t)})$ is called a \emph{generalized EM} step.
This generalized form is standard in the modern EM literature; see, for example, \citet[\S2.1]{balakrishnan_wainwright_yu_2014_em}.
\Cref{eq:em-q-function} is the operational form of the same free-energy problem introduced above, because
\begin{equation}
  \mathcal{F}(q,\theta)
  \coloneqq
  \E_q[\log p_\theta(x,Z)] + H(q)
  =
  \int q(z)\log \frac{p_\theta(x,z)}{q(z)}\,dz.
  \label{eq:em-free-energy}
\end{equation}
EM is exact coordinate ascent on this joint objective.

\begin{proposition}[EM as coordinate ascent on free energy]
\label{prop:em-free-energy-coordinate-ascent}
For any density $q$ supported where $p_\theta(x,z)>0$,
\begin{equation}
  \ell(\theta;x)
  =
  \mathcal{F}(q,\theta)
  +
  \operatorname{KL}\!\bigl(q(z)\,\|\,p_\theta(z\mid x)\bigr).
  \label{eq:em-free-energy-identity}
\end{equation}
Consequently:
\begin{enumerate}[leftmargin=*]
  \item for fixed $\theta$, the unique maximizer of $\mathcal{F}(q,\theta)$ over all densities $q$ is the exact posterior $p_\theta(z\mid x)$;
  \item exact EM alternates the exact $q$-update
  \[
    q^{(t+1)}(z)=p_{\theta^{(t)}}(z\mid x)
  \]
  with the exact $\theta$-update
  \[
    \theta^{(t+1)}\in\argmax_\theta \mathcal{F}(q^{(t+1)},\theta);
  \]
  \item if $q$ is restricted to a tractable family $\mathcal{Q}$, then maximizing $\mathcal{F}(q,\theta)$ over $q\in\mathcal{Q}$ and $\theta$ is precisely the variational-inference relaxation of EM.
\end{enumerate}
\end{proposition}
\begin{proof}
Starting from Bayes' rule,
\[
  p_\theta(z\mid x)=\frac{p_\theta(x,z)}{p_\theta(x)}.
\]
Substituting this identity into the definition of KL divergence gives
\[
  \operatorname{KL}\!\bigl(q(z)\,\|\,p_\theta(z\mid x)\bigr)
  =
  \int q(z)\log \frac{q(z)\,p_\theta(x)}{p_\theta(x,z)}\,dz.
\]
Rearranging yields \eqref{eq:em-free-energy-identity}.
For fixed $\theta$, the term \(\ell(\theta;x)\) is constant in \(q\), so maximizing \(\mathcal{F}(q,\theta)\) is equivalent to minimizing the KL divergence to the exact posterior.
That proves the first claim.

For the second claim, the exact $q$-update is just the maximizer from the first part.
Once $q^{(t+1)}$ is fixed, the entropy term $H(q^{(t+1)})$ is constant in $\theta$, so maximizing $\mathcal{F}(q^{(t+1)},\theta)$ over $\theta$ is equivalent to maximizing
\[
  \E_{q^{(t+1)}}[\log p_\theta(x,Z)]
  =
  Q(\theta\mid \theta^{(t)}),
\]
which is exactly the M-step.
The third claim is the same statement with the $q$-optimization restricted from all densities to the smaller family $\mathcal{Q}$.
\end{proof}

\begin{proposition}[EM monotonicity]
\label{prop:em-monotonicity}
Let $\ell(\theta;x)=\log p_\theta(x)$ and let $q_t(z)=p_{\theta^{(t)}}(z\mid x)$.
If $\theta^{(t+1)}$ satisfies
\[
  Q(\theta^{(t+1)}\mid \theta^{(t)}) \ge Q(\theta^{(t)}\mid \theta^{(t)}),
\]
then
\[
  \ell(\theta^{(t+1)};x) \ge \ell(\theta^{(t)};x).
\]
In particular, exact EM does not decrease the observed-data log-likelihood.
\end{proposition}
\begin{proof}
For any $\theta$,
\[
  \ell(\theta;x)
  =
  \log \int q_t(z)\,\frac{p_\theta(x,z)}{q_t(z)}\,dz.
\]
Applying Jensen's inequality to the concave logarithm gives
\[
  \ell(\theta;x)
  \ge
  \int q_t(z)\log \frac{p_\theta(x,z)}{q_t(z)}\,dz
  =
  Q(\theta\mid \theta^{(t)}) - \int q_t(z)\log q_t(z)\,dz.
\]
The second term does not depend on $\theta$.
Moreover, when $\theta=\theta^{(t)}$, Bayes' rule gives
\[
  \frac{p_{\theta^{(t)}}(x,z)}{q_t(z)}
  =
  \frac{p_{\theta^{(t)}}(x,z)}{p_{\theta^{(t)}}(z\mid x)}
  =
  p_{\theta^{(t)}}(x),
\]
which is constant in $z$, so Jensen's inequality is tight at $\theta^{(t)}$.
Therefore
\[
  \ell(\theta^{(t+1)};x) - \ell(\theta^{(t)};x)
  \ge
  Q(\theta^{(t+1)}\mid \theta^{(t)}) - Q(\theta^{(t)}\mid \theta^{(t)})
  \ge 0.
\]
\end{proof}

This proposition explains why EM is often described as a monotone ascent method.
It is important, however, not to over-interpret the result:
monotonicity guarantees improvement in likelihood, but it does \emph{not} guarantee convergence to the global maximizer.
Observed-data likelihoods for latent-variable models are usually nonconcave, so EM should be viewed as a local refinement method.

\subsubsection{Worked example: a two-coin mixture}
\label{subsubsec:em-two-coins}

Suppose coin $A$ lands heads with probability $\theta_A$ and coin $B$ lands heads with probability $\theta_B$.
For each experiment $i=1,\dots,N$:
\begin{enumerate}[leftmargin=*]
  \item choose coin $A$ with probability $\pi$ and coin $B$ with probability $1-\pi$,
  \item toss the chosen coin $m$ times,
  \item record only the number of heads $X_i\in\{0,1,\dots,m\}$.
\end{enumerate}
The latent variable is
\[
  Z_i \in \{A,B\},
\]
which records which coin was used in experiment $i$.
Conditional on $Z_i=A$, we have $X_i\sim \operatorname{Bin}(m,\theta_A)$, and conditional on $Z_i=B$, we have $X_i\sim \operatorname{Bin}(m,\theta_B)$.

If the labels $Z_1,\dots,Z_N$ were observed, estimation would be straightforward.
Indeed, the complete-data likelihood is
\begin{align}
  p(x,z\mid \pi,\theta_A,\theta_B)
  =
  \prod_{i=1}^N
  &\Bigl[
    \pi \binom{m}{x_i}\theta_A^{x_i}(1-\theta_A)^{m-x_i}
  \Bigr]^{\1\{z_i=A\}}
  \nonumber\\
  &\times
  \Bigl[
    (1-\pi)\binom{m}{x_i}\theta_B^{x_i}(1-\theta_B)^{m-x_i}
  \Bigr]^{\1\{z_i=B\}}.
  \label{eq:em-two-coins-complete}
\end{align}
The corresponding complete-data log-likelihood is
\begin{align}
  \log p(x,z\mid \pi,\theta_A,\theta_B)
  =
  \sum_{i=1}^N
  \Bigl[
    &\1\{z_i=A\}\bigl(\log \pi + x_i\log \theta_A + (m-x_i)\log(1-\theta_A)\bigr)
    \nonumber\\
    +
    &\1\{z_i=B\}\bigl(\log(1-\pi) + x_i\log \theta_B + (m-x_i)\log(1-\theta_B)\bigr)
    \nonumber\\
    +
    &\log \binom{m}{x_i}
  \Bigr].
  \label{eq:em-two-coins-complete-loglik}
\end{align}
If the indicators were known, maximizing \eqref{eq:em-two-coins-complete-loglik} would just mean taking the empirical fraction of experiments assigned to each coin and the empirical fraction of heads within each coin.

The observed-data likelihood is harder because we must sum over the unknown labels:
\begin{equation}
  p(x\mid \pi,\theta_A,\theta_B)
  =
  \prod_{i=1}^N
  \left[
    \pi \binom{m}{x_i}\theta_A^{x_i}(1-\theta_A)^{m-x_i}
    +
    (1-\pi)\binom{m}{x_i}\theta_B^{x_i}(1-\theta_B)^{m-x_i}
  \right].
  \label{eq:em-two-coins-observed-likelihood}
\end{equation}
Taking a logarithm turns the product into a sum, but each summand still contains a logarithm of a \emph{sum}, so direct maximization is awkward.

Given current estimates $(\pi^{(t)},\theta_A^{(t)},\theta_B^{(t)})$, the E-step computes the posterior probability that experiment $i$ used coin $A$:
\begin{equation}
  \gamma_{i,A}^{(t)}
  \coloneqq
  \Pr(Z_i=A\mid X_i=x_i,\pi^{(t)},\theta_A^{(t)},\theta_B^{(t)})
  =
  \frac{
    \pi^{(t)} (\theta_A^{(t)})^{x_i}(1-\theta_A^{(t)})^{m-x_i}
  }{
    \pi^{(t)} (\theta_A^{(t)})^{x_i}(1-\theta_A^{(t)})^{m-x_i}
    +
    (1-\pi^{(t)}) (\theta_B^{(t)})^{x_i}(1-\theta_B^{(t)})^{m-x_i}
  }.
  \label{eq:em-two-coins-responsibility}
\end{equation}
The binomial coefficient cancels in the ratio, so only the relative plausibility of the two coins matters.
We also write
\[
  \gamma_{i,B}^{(t)} \coloneqq 1-\gamma_{i,A}^{(t)}.
\]

Using these posterior weights in \eqref{eq:em-two-coins-complete-loglik} yields the surrogate
\begin{align}
  Q(\pi,\theta_A,\theta_B \mid \pi^{(t)},\theta_A^{(t)},\theta_B^{(t)})
  =
  \sum_{i=1}^N
  \Bigl[
    &\gamma_{i,A}^{(t)}\bigl(\log \pi + x_i\log \theta_A + (m-x_i)\log(1-\theta_A)\bigr)
    \nonumber\\
    +
    &\gamma_{i,B}^{(t)}\bigl(\log(1-\pi) + x_i\log \theta_B + (m-x_i)\log(1-\theta_B)\bigr)
  \Bigr]
  + \text{const},
  \label{eq:em-two-coins-q}
\end{align}
where the constant term does not depend on $(\pi,\theta_A,\theta_B)$.

The M-step maximizes \eqref{eq:em-two-coins-q}.
Because the parameters separate, the maximizers are available in closed form:
\begin{equation}
  \pi^{(t+1)} = \frac{1}{N}\sum_{i=1}^N \gamma_{i,A}^{(t)},
  \label{eq:em-two-coins-pi-update}
\end{equation}
\begin{equation}
  \theta_A^{(t+1)}
  =
  \frac{\sum_{i=1}^N \gamma_{i,A}^{(t)} x_i}{m\sum_{i=1}^N \gamma_{i,A}^{(t)}},
  \qquad
  \theta_B^{(t+1)}
  =
  \frac{\sum_{i=1}^N \gamma_{i,B}^{(t)} x_i}{m\sum_{i=1}^N \gamma_{i,B}^{(t)}}.
  \label{eq:em-two-coins-theta-update}
\end{equation}
These formulas have a simple interpretation:
EM replaces the unknown label of each experiment by a \emph{soft assignment}.
The updated coin-bias estimates are therefore just weighted fractions of heads, and the updated mixing weight is the average posterior probability of using coin $A$.
In the language of clustering, the E-step softly separates the experiments into ``coin $A$-like'' and ``coin $B$-like'' groups, and the M-step re-fits the two coins to those soft groups.

\subsubsection{Stable responsibilities and regularized updates}
\label{subsubsec:em-stability}

The algebra of the E-step is simple, but the computation can still become unstable when some latent states are much less plausible than others.
The general log-scale normalization rule from \Cref{subsec:stable-evaluation} is therefore part of the standard EM toolkit.
If a model has latent states $k=1,\dots,K$, one first forms log responsibilities
\[
  \ell_{ik}
  \coloneqq
  \log \pi_k + \log p(x_i\mid z_i=k,\theta_k),
\]
then normalizes them by subtracting a common scale before exponentiating:
\[
  r_{ik}
  =
  \exp\!\left(
    \ell_{ik} - \operatorname{LSE}(\ell_{i1},\dots,\ell_{iK})
  \right).
\]
Equivalently, one may subtract $\max_k \ell_{ik}$ first and normalize the resulting exponentials.
Both calculations preserve the same mathematical answer, but the log-scale version avoids underflow when one component is exponentially less plausible than another.

In models with covariance or precision updates, the M-step raises a second issue.
Weighted empirical covariance matrices can be unstable when a latent state receives very little effective mass, and the resulting estimate may be nearly singular even though the population problem is well posed.
A small covariance floor, diagonal loading term, or other mild regularization then plays the same role as damping in Newton's method: it preserves positive definiteness and prevents the next E-step from being driven by spurious tiny eigenvalues; see \Cref{subsec:geometry-control}.
This does not change the basic EM logic; it keeps the geometry of the update trustworthy while the effective sample size assigned to a component is still small.

Initialization remains the third computational safeguard.
Because EM is a local method, multiple dispersed starts are often as important as any single numerical trick.
They separate two different failure modes: unstable local arithmetic inside a run, and convergence to different basins across runs.
These same log-scale and geometry-control ideas reappear in more complicated latent-variable procedures, including hidden-state models and mixture models, where repeated normalization and covariance estimation are unavoidable.

\subsubsection{Initialization, local maxima, and the bridge to VI}
\label{subsubsec:em-local-behavior}

The two-coin example highlights why EM is different from a hard-imputation procedure.
One could try to replace each unobserved $Z_i$ by its most likely value under the current parameter and then maximize the complete-data likelihood.
That approach throws away uncertainty, especially for observations that are compatible with either latent state.
EM keeps this uncertainty in the E-step by averaging over the full conditional law \(p_{\theta^{(t)}}(z\mid x)\).

At the same time, EM remains a local method.
The likelihood ascent property from \cref{prop:em-monotonicity} says that EM moves uphill, but it does not say \emph{which} hill the algorithm will climb.
For latent-variable models with multiple basins of attraction, different initial conditions can lead to different fixed points or different local maxima.
Practitioners therefore frequently run EM from several dispersed initializations and retain the run with the largest final likelihood.

\begin{figure}[t]
  \centering
  \begin{tikzpicture}
  \begin{axis}[
    name=left,
    width=0.48\linewidth,
    height=6.0cm,
    xmin=-3, xmax=3,
    ymin=-3, ymax=3,
    axis lines=left,
    xlabel={$\theta$},
    ylabel={$M(\theta)$},
    title={$\theta^\star/\sigma=0.6$ (weaker separation)},
    title style={font=\small, yshift=-0.35em},
    clip=false,
  ]
    \addplot[stat221Blue,thick] table[row sep=\\] {
x y \\
-3.000000 -0.908632447092 \\
-2.975000 -0.908175672392 \\
-2.950000 -0.907708112179 \\
-2.925000 -0.907229431833 \\
-2.900000 -0.906739283873 \\
-2.875000 -0.906237307375 \\
-2.850000 -0.905723127352 \\
-2.825000 -0.905196354117 \\
-2.800000 -0.904656582593 \\
-2.775000 -0.904103391603 \\
-2.750000 -0.903536343115 \\
-2.725000 -0.902954981443 \\
-2.700000 -0.902358832414 \\
-2.675000 -0.90174740248 \\
-2.650000 -0.901120177787 \\
-2.625000 -0.900476623193 \\
-2.600000 -0.899816181227 \\
-2.575000 -0.899138270997 \\
-2.550000 -0.898442287026 \\
-2.525000 -0.897727598036 \\
-2.500000 -0.896993545651 \\
-2.475000 -0.89623944303 \\
-2.450000 -0.895464573424 \\
-2.425000 -0.894668188646 \\
-2.400000 -0.893849507451 \\
-2.375000 -0.893007713827 \\
-2.350000 -0.892141955175 \\
-2.325000 -0.891251340392 \\
-2.300000 -0.89033493783 \\
-2.275000 -0.889391773139 \\
-2.250000 -0.888420826976 \\
-2.225000 -0.887421032577 \\
-2.200000 -0.886391273176 \\
-2.175000 -0.88533037927 \\
-2.150000 -0.884237125712 \\
-2.125000 -0.883110228622 \\
-2.100000 -0.881948342109 \\
-2.075000 -0.88075005478 \\
-2.050000 -0.879513886032 \\
-2.025000 -0.878238282106 \\
-2.000000 -0.876921611887 \\
-1.975000 -0.875562162439 \\
-1.950000 -0.874158134241 \\
-1.925000 -0.872707636123 \\
-1.900000 -0.871208679861 \\
-1.875000 -0.869659174425 \\
-1.850000 -0.868056919841 \\
-1.825000 -0.866399600647 \\
-1.800000 -0.864684778912 \\
-1.775000 -0.862909886792 \\
-1.750000 -0.861072218575 \\
-1.725000 -0.859168922201 \\
-1.700000 -0.857196990197 \\
-1.675000 -0.855153250004 \\
-1.650000 -0.853034353638 \\
-1.625000 -0.850836766651 \\
-1.600000 -0.848556756326 \\
-1.575000 -0.846190379073 \\
-1.550000 -0.843733466952 \\
-1.525000 -0.84118161327 \\
-1.500000 -0.83853015719 \\
-1.475000 -0.83577416728 \\
-1.450000 -0.832908423933 \\
-1.425000 -0.829927400586 \\
-1.400000 -0.826825243651 \\
-1.375000 -0.823595751085 \\
-1.350000 -0.820232349514 \\
-1.325000 -0.816728069817 \\
-1.300000 -0.813075521085 \\
-1.275000 -0.809266862869 \\
-1.250000 -0.805293775623 \\
-1.225000 -0.801147429249 \\
-1.200000 -0.796818449682 \\
-1.175000 -0.792296883409 \\
-1.150000 -0.787572159879 \\
-1.125000 -0.782633051742 \\
-1.100000 -0.777467632878 \\
-1.075000 -0.772063234227 \\
-1.050000 -0.766406397422 \\
-1.025000 -0.760482826317 \\
-1.000000 -0.754277336514 \\
-0.975000 -0.747773803103 \\
-0.950000 -0.740955106877 \\
-0.925000 -0.733803079435 \\
-0.900000 -0.72629844767 \\
-0.875000 -0.71842077836 \\
-0.850000 -0.710148423741 \\
-0.825000 -0.701458469203 \\
-0.800000 -0.692326684565 \\
-0.775000 -0.682727480725 \\
-0.750000 -0.67263387393 \\
-0.725000 -0.662017460444 \\
-0.700000 -0.650848404994 \\
-0.675000 -0.639095447124 \\
-0.650000 -0.626725930439 \\
-0.625000 -0.613705860715 \\
-0.600000 -0.6 \\
-0.575000 -0.585572005116 \\
-0.550000 -0.570384620375 \\
-0.525000 -0.554399935874 \\
-0.500000 -0.537579724311 \\
-0.475000 -0.519885870784 \\
-0.450000 -0.501280911413 \\
-0.425000 -0.481728697539 \\
-0.400000 -0.461195202495 \\
-0.375000 -0.43964948694 \\
-0.350000 -0.417064836032 \\
-0.325000 -0.393420076445 \\
-0.300000 -0.368701072508 \\
-0.275000 -0.342902387595 \\
-0.250000 -0.316029078272 \\
-0.225000 -0.288098564002 \\
-0.200000 -0.259142484414 \\
-0.175000 -0.229208420717 \\
-0.150000 -0.198361321295 \\
-0.125000 -0.166684440279 \\
-0.100000 -0.134279581613 \\
-0.075000 -0.101266451527 \\
-0.050000 -0.0677809710687 \\
-0.025000 -0.0339724931235 \\
0.000000 -4.40420811967e-18 \\
0.025000 0.0339724931235 \\
0.050000 0.0677809710687 \\
0.075000 0.101266451527 \\
0.100000 0.134279581613 \\
0.125000 0.166684440279 \\
0.150000 0.198361321295 \\
0.175000 0.229208420717 \\
0.200000 0.259142484414 \\
0.225000 0.288098564002 \\
0.250000 0.316029078272 \\
0.275000 0.342902387595 \\
0.300000 0.368701072508 \\
0.325000 0.393420076445 \\
0.350000 0.417064836032 \\
0.375000 0.43964948694 \\
0.400000 0.461195202495 \\
0.425000 0.481728697539 \\
0.450000 0.501280911413 \\
0.475000 0.519885870784 \\
0.500000 0.537579724311 \\
0.525000 0.554399935874 \\
0.550000 0.570384620375 \\
0.575000 0.585572005116 \\
0.600000 0.6 \\
0.625000 0.613705860715 \\
0.650000 0.626725930439 \\
0.675000 0.639095447124 \\
0.700000 0.650848404994 \\
0.725000 0.662017460444 \\
0.750000 0.67263387393 \\
0.775000 0.682727480725 \\
0.800000 0.692326684565 \\
0.825000 0.701458469203 \\
0.850000 0.710148423741 \\
0.875000 0.71842077836 \\
0.900000 0.72629844767 \\
0.925000 0.733803079435 \\
0.950000 0.740955106877 \\
0.975000 0.747773803103 \\
1.000000 0.754277336514 \\
1.025000 0.760482826317 \\
1.050000 0.766406397422 \\
1.075000 0.772063234227 \\
1.100000 0.777467632878 \\
1.125000 0.782633051742 \\
1.150000 0.787572159879 \\
1.175000 0.792296883409 \\
1.200000 0.796818449682 \\
1.225000 0.801147429249 \\
1.250000 0.805293775623 \\
1.275000 0.809266862869 \\
1.300000 0.813075521085 \\
1.325000 0.816728069817 \\
1.350000 0.820232349514 \\
1.375000 0.823595751085 \\
1.400000 0.826825243651 \\
1.425000 0.829927400586 \\
1.450000 0.832908423933 \\
1.475000 0.83577416728 \\
1.500000 0.83853015719 \\
1.525000 0.84118161327 \\
1.550000 0.843733466952 \\
1.575000 0.846190379073 \\
1.600000 0.848556756326 \\
1.625000 0.850836766651 \\
1.650000 0.853034353638 \\
1.675000 0.855153250004 \\
1.700000 0.857196990197 \\
1.725000 0.859168922201 \\
1.750000 0.861072218575 \\
1.775000 0.862909886792 \\
1.800000 0.864684778912 \\
1.825000 0.866399600647 \\
1.850000 0.868056919841 \\
1.875000 0.869659174425 \\
1.900000 0.871208679861 \\
1.925000 0.872707636123 \\
1.950000 0.874158134241 \\
1.975000 0.875562162439 \\
2.000000 0.876921611887 \\
2.025000 0.878238282106 \\
2.050000 0.879513886032 \\
2.075000 0.88075005478 \\
2.100000 0.881948342109 \\
2.125000 0.883110228622 \\
2.150000 0.884237125712 \\
2.175000 0.88533037927 \\
2.200000 0.886391273176 \\
2.225000 0.887421032577 \\
2.250000 0.888420826976 \\
2.275000 0.889391773139 \\
2.300000 0.89033493783 \\
2.325000 0.891251340392 \\
2.350000 0.892141955175 \\
2.375000 0.893007713827 \\
2.400000 0.893849507451 \\
2.425000 0.894668188646 \\
2.450000 0.895464573424 \\
2.475000 0.89623944303 \\
2.500000 0.896993545651 \\
2.525000 0.897727598036 \\
2.550000 0.898442287026 \\
2.575000 0.899138270997 \\
2.600000 0.899816181227 \\
2.625000 0.900476623193 \\
2.650000 0.901120177787 \\
2.675000 0.90174740248 \\
2.700000 0.902358832414 \\
2.725000 0.902954981443 \\
2.750000 0.903536343115 \\
2.775000 0.904103391603 \\
2.800000 0.904656582593 \\
2.825000 0.905196354117 \\
2.850000 0.905723127352 \\
2.875000 0.906237307375 \\
2.900000 0.906739283873 \\
2.925000 0.907229431833 \\
2.950000 0.907708112179 \\
2.975000 0.908175672392 \\
3.000000 0.908632447092 \\
    };
    \addplot[stat221Gray,densely dashed,thick,opacity=0.9] table[row sep=\\] {
x y \\
-3.000000 -3 \\
-2.975000 -2.975 \\
-2.950000 -2.95 \\
-2.925000 -2.925 \\
-2.900000 -2.9 \\
-2.875000 -2.875 \\
-2.850000 -2.85 \\
-2.825000 -2.825 \\
-2.800000 -2.8 \\
-2.775000 -2.775 \\
-2.750000 -2.75 \\
-2.725000 -2.725 \\
-2.700000 -2.7 \\
-2.675000 -2.675 \\
-2.650000 -2.65 \\
-2.625000 -2.625 \\
-2.600000 -2.6 \\
-2.575000 -2.575 \\
-2.550000 -2.55 \\
-2.525000 -2.525 \\
-2.500000 -2.5 \\
-2.475000 -2.475 \\
-2.450000 -2.45 \\
-2.425000 -2.425 \\
-2.400000 -2.4 \\
-2.375000 -2.375 \\
-2.350000 -2.35 \\
-2.325000 -2.325 \\
-2.300000 -2.3 \\
-2.275000 -2.275 \\
-2.250000 -2.25 \\
-2.225000 -2.225 \\
-2.200000 -2.2 \\
-2.175000 -2.175 \\
-2.150000 -2.15 \\
-2.125000 -2.125 \\
-2.100000 -2.1 \\
-2.075000 -2.075 \\
-2.050000 -2.05 \\
-2.025000 -2.025 \\
-2.000000 -2 \\
-1.975000 -1.975 \\
-1.950000 -1.95 \\
-1.925000 -1.925 \\
-1.900000 -1.9 \\
-1.875000 -1.875 \\
-1.850000 -1.85 \\
-1.825000 -1.825 \\
-1.800000 -1.8 \\
-1.775000 -1.775 \\
-1.750000 -1.75 \\
-1.725000 -1.725 \\
-1.700000 -1.7 \\
-1.675000 -1.675 \\
-1.650000 -1.65 \\
-1.625000 -1.625 \\
-1.600000 -1.6 \\
-1.575000 -1.575 \\
-1.550000 -1.55 \\
-1.525000 -1.525 \\
-1.500000 -1.5 \\
-1.475000 -1.475 \\
-1.450000 -1.45 \\
-1.425000 -1.425 \\
-1.400000 -1.4 \\
-1.375000 -1.375 \\
-1.350000 -1.35 \\
-1.325000 -1.325 \\
-1.300000 -1.3 \\
-1.275000 -1.275 \\
-1.250000 -1.25 \\
-1.225000 -1.225 \\
-1.200000 -1.2 \\
-1.175000 -1.175 \\
-1.150000 -1.15 \\
-1.125000 -1.125 \\
-1.100000 -1.1 \\
-1.075000 -1.075 \\
-1.050000 -1.05 \\
-1.025000 -1.025 \\
-1.000000 -1 \\
-0.975000 -0.975 \\
-0.950000 -0.95 \\
-0.925000 -0.925 \\
-0.900000 -0.9 \\
-0.875000 -0.875 \\
-0.850000 -0.85 \\
-0.825000 -0.825 \\
-0.800000 -0.8 \\
-0.775000 -0.775 \\
-0.750000 -0.75 \\
-0.725000 -0.725 \\
-0.700000 -0.7 \\
-0.675000 -0.675 \\
-0.650000 -0.65 \\
-0.625000 -0.625 \\
-0.600000 -0.6 \\
-0.575000 -0.575 \\
-0.550000 -0.55 \\
-0.525000 -0.525 \\
-0.500000 -0.5 \\
-0.475000 -0.475 \\
-0.450000 -0.45 \\
-0.425000 -0.425 \\
-0.400000 -0.4 \\
-0.375000 -0.375 \\
-0.350000 -0.35 \\
-0.325000 -0.325 \\
-0.300000 -0.3 \\
-0.275000 -0.275 \\
-0.250000 -0.25 \\
-0.225000 -0.225 \\
-0.200000 -0.2 \\
-0.175000 -0.175 \\
-0.150000 -0.15 \\
-0.125000 -0.125 \\
-0.100000 -0.1 \\
-0.075000 -0.075 \\
-0.050000 -0.05 \\
-0.025000 -0.025 \\
0.000000 0 \\
0.025000 0.025 \\
0.050000 0.05 \\
0.075000 0.075 \\
0.100000 0.1 \\
0.125000 0.125 \\
0.150000 0.15 \\
0.175000 0.175 \\
0.200000 0.2 \\
0.225000 0.225 \\
0.250000 0.25 \\
0.275000 0.275 \\
0.300000 0.3 \\
0.325000 0.325 \\
0.350000 0.35 \\
0.375000 0.375 \\
0.400000 0.4 \\
0.425000 0.425 \\
0.450000 0.45 \\
0.475000 0.475 \\
0.500000 0.5 \\
0.525000 0.525 \\
0.550000 0.55 \\
0.575000 0.575 \\
0.600000 0.6 \\
0.625000 0.625 \\
0.650000 0.65 \\
0.675000 0.675 \\
0.700000 0.7 \\
0.725000 0.725 \\
0.750000 0.75 \\
0.775000 0.775 \\
0.800000 0.8 \\
0.825000 0.825 \\
0.850000 0.85 \\
0.875000 0.875 \\
0.900000 0.9 \\
0.925000 0.925 \\
0.950000 0.95 \\
0.975000 0.975 \\
1.000000 1 \\
1.025000 1.025 \\
1.050000 1.05 \\
1.075000 1.075 \\
1.100000 1.1 \\
1.125000 1.125 \\
1.150000 1.15 \\
1.175000 1.175 \\
1.200000 1.2 \\
1.225000 1.225 \\
1.250000 1.25 \\
1.275000 1.275 \\
1.300000 1.3 \\
1.325000 1.325 \\
1.350000 1.35 \\
1.375000 1.375 \\
1.400000 1.4 \\
1.425000 1.425 \\
1.450000 1.45 \\
1.475000 1.475 \\
1.500000 1.5 \\
1.525000 1.525 \\
1.550000 1.55 \\
1.575000 1.575 \\
1.600000 1.6 \\
1.625000 1.625 \\
1.650000 1.65 \\
1.675000 1.675 \\
1.700000 1.7 \\
1.725000 1.725 \\
1.750000 1.75 \\
1.775000 1.775 \\
1.800000 1.8 \\
1.825000 1.825 \\
1.850000 1.85 \\
1.875000 1.875 \\
1.900000 1.9 \\
1.925000 1.925 \\
1.950000 1.95 \\
1.975000 1.975 \\
2.000000 2 \\
2.025000 2.025 \\
2.050000 2.05 \\
2.075000 2.075 \\
2.100000 2.1 \\
2.125000 2.125 \\
2.150000 2.15 \\
2.175000 2.175 \\
2.200000 2.2 \\
2.225000 2.225 \\
2.250000 2.25 \\
2.275000 2.275 \\
2.300000 2.3 \\
2.325000 2.325 \\
2.350000 2.35 \\
2.375000 2.375 \\
2.400000 2.4 \\
2.425000 2.425 \\
2.450000 2.45 \\
2.475000 2.475 \\
2.500000 2.5 \\
2.525000 2.525 \\
2.550000 2.55 \\
2.575000 2.575 \\
2.600000 2.6 \\
2.625000 2.625 \\
2.650000 2.65 \\
2.675000 2.675 \\
2.700000 2.7 \\
2.725000 2.725 \\
2.750000 2.75 \\
2.775000 2.775 \\
2.800000 2.8 \\
2.825000 2.825 \\
2.850000 2.85 \\
2.875000 2.875 \\
2.900000 2.9 \\
2.925000 2.925 \\
2.950000 2.95 \\
2.975000 2.975 \\
3.000000 3 \\
    };

    \addplot[stat221Orange,thick,mark=*,mark size=0.7pt,opacity=0.85] table[row sep=\\] {
x y \\
0.200000 0.2 \\
0.200000 0.259142484414 \\
0.259142 0.259142484414 \\
0.259142 0.325980454969 \\
0.325980 0.325980454969 \\
0.325980 0.394367553471 \\
0.394368 0.394367553471 \\
0.394368 0.456430362406 \\
0.456430 0.456430362406 \\
0.456430 0.506155448593 \\
0.506155 0.506155448593 \\
0.506155 0.541800771088 \\
0.541801 0.541800771088 \\
    };

    \addplot[only marks,mark=o,mark size=1.6pt,stat221Gray,fill=stat221Gray] coordinates {(0,0)};
    \addplot[only marks,mark=x,mark size=2.7pt,stat221Purple,very thick] coordinates {(-0.600000, -0.600000)};
    \addplot[only marks,mark=x,mark size=2.7pt,stat221Red,very thick] coordinates {(0.600000, 0.600000)};

    \node[anchor=south west] at (axis cs:0.600000,0.600000) {$\theta^\star$};
    \node[anchor=south east] at (axis cs:-0.600000,-0.600000) {$-\theta^\star$};
    \node[anchor=south west] at (axis cs:0,0) {$0$};

    \node[stat221Blue,anchor=south west,font=\small] at (axis cs:-2.85,2.55) {$M(\theta)$};
    \node[stat221Gray,anchor=south west,font=\small] at (axis cs:-2.85,2.25) {$\theta$};
    \node[stat221Orange,anchor=south west,font=\small] at (axis cs:-2.85,1.95) {cobweb from $\theta^0=0.2$};
    \node[anchor=south west,font=\small] at (axis cs:-2.85,1.65) {$|M'(\theta^\star)|\approx 0.562$};
  \end{axis}

  \begin{axis}[
    at={(left.north east)}, anchor=north west,
    xshift=0.45cm,
    width=0.48\linewidth,
    height=6.0cm,
    xmin=-3, xmax=3,
    ymin=-3, ymax=3,
    axis lines=left,
    xlabel={$\theta$},
    ylabel={},
    title={$\theta^\star/\sigma=2.0$ (stronger separation)},
    title style={font=\small, yshift=-0.35em},
    clip=false,
    yticklabels={},
  ]
    \addplot[stat221Blue,thick] table[row sep=\\] {
x y \\
-3.000000 -2.01074543078 \\
-2.975000 -2.01061987459 \\
-2.950000 -2.01049048553 \\
-2.925000 -2.01035710807 \\
-2.900000 -2.01021957892 \\
-2.875000 -2.01007772657 \\
-2.850000 -2.00993137075 \\
-2.825000 -2.00978032198 \\
-2.800000 -2.00962438098 \\
-2.775000 -2.00946333804 \\
-2.750000 -2.00929697244 \\
-2.725000 -2.00912505172 \\
-2.700000 -2.00894733097 \\
-2.675000 -2.00876355204 \\
-2.650000 -2.00857344269 \\
-2.625000 -2.0083767157 \\
-2.600000 -2.00817306789 \\
-2.575000 -2.00796217911 \\
-2.550000 -2.00774371107 \\
-2.525000 -2.0075173062 \\
-2.500000 -2.00728258631 \\
-2.475000 -2.00703915123 \\
-2.450000 -2.00678657732 \\
-2.425000 -2.00652441583 \\
-2.400000 -2.00625219119 \\
-2.375000 -2.00596939912 \\
-2.350000 -2.00567550463 \\
-2.325000 -2.00536993984 \\
-2.300000 -2.00505210158 \\
-2.275000 -2.0047213489 \\
-2.250000 -2.00437700029 \\
-2.225000 -2.00401833068 \\
-2.200000 -2.00364456824 \\
-2.175000 -2.00325489085 \\
-2.150000 -2.00284842232 \\
-2.125000 -2.00242422828 \\
-2.100000 -2.00198131166 \\
-2.075000 -2.00151860788 \\
-2.050000 -2.00103497953 \\
-2.025000 -2.00052921062 \\
-2.000000 -2.00000000033 \\
-1.975000 -1.99944595616 \\
-1.950000 -1.99886558653 \\
-1.925000 -1.9982572926 \\
-1.900000 -1.99761935948 \\
-1.875000 -1.99694994645 \\
-1.850000 -1.99624707648 \\
-1.825000 -1.99550862454 \\
-1.800000 -1.99473230499 \\
-1.775000 -1.99391565764 \\
-1.750000 -1.99305603251 \\
-1.725000 -1.99215057315 \\
-1.700000 -1.99119619819 \\
-1.675000 -1.99018958127 \\
-1.650000 -1.98912712875 \\
-1.625000 -1.98800495536 \\
-1.600000 -1.9868188573 \\
-1.575000 -1.98556428257 \\
-1.550000 -1.98423629834 \\
-1.525000 -1.98282955485 \\
-1.500000 -1.9813382456 \\
-1.475000 -1.97975606334 \\
-1.450000 -1.97807615137 \\
-1.425000 -1.97629104969 \\
-1.400000 -1.97439263535 \\
-1.375000 -1.97237205633 \\
-1.350000 -1.97021965824 \\
-1.325000 -1.96792490303 \\
-1.300000 -1.96547627869 \\
-1.275000 -1.96286119896 \\
-1.250000 -1.96006589192 \\
-1.225000 -1.95707527605 \\
-1.200000 -1.9538728224 \\
-1.175000 -1.9504404012 \\
-1.150000 -1.94675811114 \\
-1.125000 -1.94280408921 \\
-1.100000 -1.93855429895 \\
-1.075000 -1.93398229465 \\
-1.050000 -1.92905895869 \\
-1.025000 -1.92375220897 \\
-1.000000 -1.9180266733 \\
-0.975000 -1.91184332696 \\
-0.950000 -1.90515908967 \\
-0.925000 -1.89792637777 \\
-0.900000 -1.89009260715 \\
-0.875000 -1.88159964246 \\
-0.850000 -1.87238318777 \\
-0.825000 -1.8623721142 \\
-0.800000 -1.85148772 \\
-0.775000 -1.83964291948 \\
-0.750000 -1.82674135777 \\
-0.725000 -1.81267645029 \\
-0.700000 -1.79733034788 \\
-0.675000 -1.78057283207 \\
-0.650000 -1.76226014974 \\
-0.625000 -1.74223380268 \\
-0.600000 -1.72031931697 \\
-0.575000 -1.69632502861 \\
-0.550000 -1.670040938 \\
-0.525000 -1.64123770626 \\
-0.500000 -1.60966589311 \\
-0.475000 -1.57505556962 \\
-0.450000 -1.53711648088 \\
-0.425000 -1.49553898428 \\
-0.400000 -1.44999604862 \\
-0.375000 -1.40014666593 \\
-0.350000 -1.34564109786 \\
-0.325000 -1.28612844298 \\
-0.300000 -1.22126705591 \\
-0.275000 -1.15073834966 \\
-0.250000 -1.07426443318 \\
-0.225000 -0.9916298287 \\
-0.200000 -0.902707117245 \\
-0.175000 -0.807485718017 \\
-0.150000 -0.706102081025 \\
-0.125000 -0.598868397852 \\
-0.100000 -0.48629567527 \\
-0.075000 -0.369106023452 \\
-0.050000 -0.248228826879 \\
-0.025000 -0.124776689047 \\
0.000000 -5.62759926402e-18 \\
0.025000 0.124776689047 \\
0.050000 0.248228826879 \\
0.075000 0.369106023452 \\
0.100000 0.48629567527 \\
0.125000 0.598868397852 \\
0.150000 0.706102081025 \\
0.175000 0.807485718017 \\
0.200000 0.902707117245 \\
0.225000 0.9916298287 \\
0.250000 1.07426443318 \\
0.275000 1.15073834966 \\
0.300000 1.22126705591 \\
0.325000 1.28612844298 \\
0.350000 1.34564109786 \\
0.375000 1.40014666593 \\
0.400000 1.44999604862 \\
0.425000 1.49553898428 \\
0.450000 1.53711648088 \\
0.475000 1.57505556962 \\
0.500000 1.60966589311 \\
0.525000 1.64123770626 \\
0.550000 1.670040938 \\
0.575000 1.69632502861 \\
0.600000 1.72031931697 \\
0.625000 1.74223380268 \\
0.650000 1.76226014974 \\
0.675000 1.78057283207 \\
0.700000 1.79733034788 \\
0.725000 1.81267645029 \\
0.750000 1.82674135777 \\
0.775000 1.83964291948 \\
0.800000 1.85148772 \\
0.825000 1.8623721142 \\
0.850000 1.87238318777 \\
0.875000 1.88159964246 \\
0.900000 1.89009260715 \\
0.925000 1.89792637777 \\
0.950000 1.90515908967 \\
0.975000 1.91184332696 \\
1.000000 1.9180266733 \\
1.025000 1.92375220897 \\
1.050000 1.92905895869 \\
1.075000 1.93398229465 \\
1.100000 1.93855429895 \\
1.125000 1.94280408921 \\
1.150000 1.94675811114 \\
1.175000 1.9504404012 \\
1.200000 1.9538728224 \\
1.225000 1.95707527605 \\
1.250000 1.96006589192 \\
1.275000 1.96286119896 \\
1.300000 1.96547627869 \\
1.325000 1.96792490303 \\
1.350000 1.97021965824 \\
1.375000 1.97237205633 \\
1.400000 1.97439263535 \\
1.425000 1.97629104969 \\
1.450000 1.97807615137 \\
1.475000 1.97975606334 \\
1.500000 1.9813382456 \\
1.525000 1.98282955485 \\
1.550000 1.98423629834 \\
1.575000 1.98556428257 \\
1.600000 1.9868188573 \\
1.625000 1.98800495536 \\
1.650000 1.98912712875 \\
1.675000 1.99018958127 \\
1.700000 1.99119619819 \\
1.725000 1.99215057315 \\
1.750000 1.99305603251 \\
1.775000 1.99391565764 \\
1.800000 1.99473230499 \\
1.825000 1.99550862454 \\
1.850000 1.99624707648 \\
1.875000 1.99694994645 \\
1.900000 1.99761935948 \\
1.925000 1.9982572926 \\
1.950000 1.99886558653 \\
1.975000 1.99944595616 \\
2.000000 2.00000000033 \\
2.025000 2.00052921062 \\
2.050000 2.00103497953 \\
2.075000 2.00151860788 \\
2.100000 2.00198131166 \\
2.125000 2.00242422828 \\
2.150000 2.00284842232 \\
2.175000 2.00325489085 \\
2.200000 2.00364456824 \\
2.225000 2.00401833068 \\
2.250000 2.00437700029 \\
2.275000 2.0047213489 \\
2.300000 2.00505210158 \\
2.325000 2.00536993984 \\
2.350000 2.00567550463 \\
2.375000 2.00596939912 \\
2.400000 2.00625219119 \\
2.425000 2.00652441583 \\
2.450000 2.00678657732 \\
2.475000 2.00703915123 \\
2.500000 2.00728258631 \\
2.525000 2.0075173062 \\
2.550000 2.00774371107 \\
2.575000 2.00796217911 \\
2.600000 2.00817306789 \\
2.625000 2.0083767157 \\
2.650000 2.00857344269 \\
2.675000 2.00876355204 \\
2.700000 2.00894733097 \\
2.725000 2.00912505172 \\
2.750000 2.00929697244 \\
2.775000 2.00946333804 \\
2.800000 2.00962438098 \\
2.825000 2.00978032198 \\
2.850000 2.00993137075 \\
2.875000 2.01007772657 \\
2.900000 2.01021957892 \\
2.925000 2.01035710807 \\
2.950000 2.01049048553 \\
2.975000 2.01061987459 \\
3.000000 2.01074543078 \\
    };
    \addplot[stat221Gray,densely dashed,thick,opacity=0.9] table[row sep=\\] {
x y \\
-3.000000 -3 \\
-2.975000 -2.975 \\
-2.950000 -2.95 \\
-2.925000 -2.925 \\
-2.900000 -2.9 \\
-2.875000 -2.875 \\
-2.850000 -2.85 \\
-2.825000 -2.825 \\
-2.800000 -2.8 \\
-2.775000 -2.775 \\
-2.750000 -2.75 \\
-2.725000 -2.725 \\
-2.700000 -2.7 \\
-2.675000 -2.675 \\
-2.650000 -2.65 \\
-2.625000 -2.625 \\
-2.600000 -2.6 \\
-2.575000 -2.575 \\
-2.550000 -2.55 \\
-2.525000 -2.525 \\
-2.500000 -2.5 \\
-2.475000 -2.475 \\
-2.450000 -2.45 \\
-2.425000 -2.425 \\
-2.400000 -2.4 \\
-2.375000 -2.375 \\
-2.350000 -2.35 \\
-2.325000 -2.325 \\
-2.300000 -2.3 \\
-2.275000 -2.275 \\
-2.250000 -2.25 \\
-2.225000 -2.225 \\
-2.200000 -2.2 \\
-2.175000 -2.175 \\
-2.150000 -2.15 \\
-2.125000 -2.125 \\
-2.100000 -2.1 \\
-2.075000 -2.075 \\
-2.050000 -2.05 \\
-2.025000 -2.025 \\
-2.000000 -2 \\
-1.975000 -1.975 \\
-1.950000 -1.95 \\
-1.925000 -1.925 \\
-1.900000 -1.9 \\
-1.875000 -1.875 \\
-1.850000 -1.85 \\
-1.825000 -1.825 \\
-1.800000 -1.8 \\
-1.775000 -1.775 \\
-1.750000 -1.75 \\
-1.725000 -1.725 \\
-1.700000 -1.7 \\
-1.675000 -1.675 \\
-1.650000 -1.65 \\
-1.625000 -1.625 \\
-1.600000 -1.6 \\
-1.575000 -1.575 \\
-1.550000 -1.55 \\
-1.525000 -1.525 \\
-1.500000 -1.5 \\
-1.475000 -1.475 \\
-1.450000 -1.45 \\
-1.425000 -1.425 \\
-1.400000 -1.4 \\
-1.375000 -1.375 \\
-1.350000 -1.35 \\
-1.325000 -1.325 \\
-1.300000 -1.3 \\
-1.275000 -1.275 \\
-1.250000 -1.25 \\
-1.225000 -1.225 \\
-1.200000 -1.2 \\
-1.175000 -1.175 \\
-1.150000 -1.15 \\
-1.125000 -1.125 \\
-1.100000 -1.1 \\
-1.075000 -1.075 \\
-1.050000 -1.05 \\
-1.025000 -1.025 \\
-1.000000 -1 \\
-0.975000 -0.975 \\
-0.950000 -0.95 \\
-0.925000 -0.925 \\
-0.900000 -0.9 \\
-0.875000 -0.875 \\
-0.850000 -0.85 \\
-0.825000 -0.825 \\
-0.800000 -0.8 \\
-0.775000 -0.775 \\
-0.750000 -0.75 \\
-0.725000 -0.725 \\
-0.700000 -0.7 \\
-0.675000 -0.675 \\
-0.650000 -0.65 \\
-0.625000 -0.625 \\
-0.600000 -0.6 \\
-0.575000 -0.575 \\
-0.550000 -0.55 \\
-0.525000 -0.525 \\
-0.500000 -0.5 \\
-0.475000 -0.475 \\
-0.450000 -0.45 \\
-0.425000 -0.425 \\
-0.400000 -0.4 \\
-0.375000 -0.375 \\
-0.350000 -0.35 \\
-0.325000 -0.325 \\
-0.300000 -0.3 \\
-0.275000 -0.275 \\
-0.250000 -0.25 \\
-0.225000 -0.225 \\
-0.200000 -0.2 \\
-0.175000 -0.175 \\
-0.150000 -0.15 \\
-0.125000 -0.125 \\
-0.100000 -0.1 \\
-0.075000 -0.075 \\
-0.050000 -0.05 \\
-0.025000 -0.025 \\
0.000000 0 \\
0.025000 0.025 \\
0.050000 0.05 \\
0.075000 0.075 \\
0.100000 0.1 \\
0.125000 0.125 \\
0.150000 0.15 \\
0.175000 0.175 \\
0.200000 0.2 \\
0.225000 0.225 \\
0.250000 0.25 \\
0.275000 0.275 \\
0.300000 0.3 \\
0.325000 0.325 \\
0.350000 0.35 \\
0.375000 0.375 \\
0.400000 0.4 \\
0.425000 0.425 \\
0.450000 0.45 \\
0.475000 0.475 \\
0.500000 0.5 \\
0.525000 0.525 \\
0.550000 0.55 \\
0.575000 0.575 \\
0.600000 0.6 \\
0.625000 0.625 \\
0.650000 0.65 \\
0.675000 0.675 \\
0.700000 0.7 \\
0.725000 0.725 \\
0.750000 0.75 \\
0.775000 0.775 \\
0.800000 0.8 \\
0.825000 0.825 \\
0.850000 0.85 \\
0.875000 0.875 \\
0.900000 0.9 \\
0.925000 0.925 \\
0.950000 0.95 \\
0.975000 0.975 \\
1.000000 1 \\
1.025000 1.025 \\
1.050000 1.05 \\
1.075000 1.075 \\
1.100000 1.1 \\
1.125000 1.125 \\
1.150000 1.15 \\
1.175000 1.175 \\
1.200000 1.2 \\
1.225000 1.225 \\
1.250000 1.25 \\
1.275000 1.275 \\
1.300000 1.3 \\
1.325000 1.325 \\
1.350000 1.35 \\
1.375000 1.375 \\
1.400000 1.4 \\
1.425000 1.425 \\
1.450000 1.45 \\
1.475000 1.475 \\
1.500000 1.5 \\
1.525000 1.525 \\
1.550000 1.55 \\
1.575000 1.575 \\
1.600000 1.6 \\
1.625000 1.625 \\
1.650000 1.65 \\
1.675000 1.675 \\
1.700000 1.7 \\
1.725000 1.725 \\
1.750000 1.75 \\
1.775000 1.775 \\
1.800000 1.8 \\
1.825000 1.825 \\
1.850000 1.85 \\
1.875000 1.875 \\
1.900000 1.9 \\
1.925000 1.925 \\
1.950000 1.95 \\
1.975000 1.975 \\
2.000000 2 \\
2.025000 2.025 \\
2.050000 2.05 \\
2.075000 2.075 \\
2.100000 2.1 \\
2.125000 2.125 \\
2.150000 2.15 \\
2.175000 2.175 \\
2.200000 2.2 \\
2.225000 2.225 \\
2.250000 2.25 \\
2.275000 2.275 \\
2.300000 2.3 \\
2.325000 2.325 \\
2.350000 2.35 \\
2.375000 2.375 \\
2.400000 2.4 \\
2.425000 2.425 \\
2.450000 2.45 \\
2.475000 2.475 \\
2.500000 2.5 \\
2.525000 2.525 \\
2.550000 2.55 \\
2.575000 2.575 \\
2.600000 2.6 \\
2.625000 2.625 \\
2.650000 2.65 \\
2.675000 2.675 \\
2.700000 2.7 \\
2.725000 2.725 \\
2.750000 2.75 \\
2.775000 2.775 \\
2.800000 2.8 \\
2.825000 2.825 \\
2.850000 2.85 \\
2.875000 2.875 \\
2.900000 2.9 \\
2.925000 2.925 \\
2.950000 2.95 \\
2.975000 2.975 \\
3.000000 3 \\
    };

    \addplot[stat221Orange,thick,mark=*,mark size=0.7pt,opacity=0.85] table[row sep=\\] {
x y \\
0.200000 0.2 \\
0.200000 0.902707117245 \\
0.902707 0.902707117245 \\
0.902707 1.89097162358 \\
1.890972 1.89097162358 \\
1.890972 1.99738134287 \\
1.997381 1.99738134287 \\
1.997381 1.99994315557 \\
1.999943 1.99994315557 \\
1.999943 1.99999876927 \\
1.999999 1.99999876927 \\
1.999999 1.99999997367 \\
2.000000 1.99999997367 \\
    };

    \addplot[only marks,mark=o,mark size=1.6pt,stat221Gray,fill=stat221Gray] coordinates {(0,0)};
    \addplot[only marks,mark=x,mark size=2.7pt,stat221Purple,very thick] coordinates {(-2.000000, -2.000000)};
    \addplot[only marks,mark=x,mark size=2.7pt,stat221Red,very thick] coordinates {(2.000000, 2.000000)};

    \node[anchor=south west] at (axis cs:2.000000,2.000000) {$\theta^\star$};
    \node[anchor=south east] at (axis cs:-2.000000,-2.000000) {$-\theta^\star$};
    \node[anchor=south west] at (axis cs:0,0) {$0$};
    \node[anchor=south west,font=\small] at (axis cs:-2.85,1.65) {$|M'(\theta^\star)|\approx 0.022$};
  \end{axis}
\end{tikzpicture}

  \caption{A one-dimensional population EM map for a symmetric Gaussian mixture.
  The fixed points are the intersections with the diagonal; the filled points are stable and the open point at the origin is unstable.
  The cobweb paths, shown here from the common initialization $\theta^{(0)}=0.2$, show how the limiting point depends on the initialization.
  Weaker separation (left) leaves the local slope near the stable fixed point around $0.56$, so refinement is slower; stronger separation (right) drives that slope down to about $0.02$, producing much sharper local contraction.
  This basin picture is why EM is best understood as a local refinement method rather than a globally convergent optimizer.}
  \label{fig:em-operator-fixed-points}
\end{figure}

The two panels of \Cref{fig:em-operator-fixed-points} sharpen two central intuitions.
First, EM is basin-based: the initialization determines which stable fixed point the iterates approach.
Second, once EM is inside the right basin, the local slope of the map controls the speed of refinement.
When the latent classes are weakly separated, the map lies closer to the diagonal near the target and convergence is slower; when the classes are better separated, the map bends more sharply and the iterates lock in much faster.
For Gaussian mixtures, these basin and contraction phenomena are formalized in both local and global analyses: \citet[Cor.~1--2]{balakrishnan_wainwright_yu_2014_em} prove local contractivity inside a basin around the truth, while \citet[Thm.~1]{daskalakis_tzamos_zampetakis_2016_ten_steps_em} and \citet[Thms.~3--4]{xu_hsu_maleki_2016_global_em} analyze the fixed-point geometry and broader convergence picture.

There is also a useful comparison with Gibbs sampling and with the coordinate-ascent variational updates that appear later in the course.
All three methods exploit latent-variable structure by alternating conditional updates, but they use those conditionals differently:
\[
  \text{Gibbs: sample } z \sim p_\theta(z\mid x), \qquad
  \text{EM: set } q(z)=p_\theta(z\mid x) \text{ and maximize in } \theta,
\]
\[
  \text{CAVI: update } q_j(z_j) \propto \exp\!\bigl(\E_{q_{-j}}[\log p_\theta(x,z)]\bigr)
  \text{ inside a restricted family.}
\]
Gibbs alternates conditional sampling, EM alternates exact conditional expectation and maximization, and CAVI keeps the same free-energy logic while restricting the latent law to a tractable family.
This is another instance of the augmented-state viewpoint used throughout this unit: introducing extra probabilistic structure can make it possible to explore or optimize a landscape in ways that purely local deterministic updates cannot.

The important conceptual handoff is that EM and VI optimize the same free-energy identity.
Exact EM uses the unrestricted update
\[
  q^{(t+1)}(z)=p_{\theta^{(t)}}(z\mid x),
\]
whereas variational inference keeps the model $p_\theta(x,z)$ fixed and replaces that exact posterior by the best approximation inside a tractable family $\mathcal{Q}$:
\[
  q^{(t+1)} \in \argmax_{q\in\mathcal{Q}} \mathcal{F}(q,\theta^{(t)}).
\]
The next unit adopts this $(x,z,\theta)$ notation throughout and rewrites the same free-energy objective as the evidence lower bound in \Cref{subsec:vi-elbo}.
